% ============================================
% Chapter 2: Understanding the Problem and the Data
% ============================================

\chapter{Understanding the Problem and the Data}\label{ch:understanding}

\section{Introduction}

This chapter translates the state of the art reviewed in Chapter~\ref{ch:state_of_art} into the concrete problem that AEGIS Entropy addresses. It follows the CPS (Context--Problematic--Solution) design methodology (Section~\ref{sec:cps}) required by the project pedagogy, reports the findings of four user interviews (Section~\ref{sec:interviews}), describes the data sources and measurement instruments (Section~\ref{sec:instruments}), and presents the data dictionary that drives the machine learning pipeline (Section~\ref{sec:data_dictionary}).

\section{Project Context --- CPS Analysis}\label{sec:cps}

\subsection{Context}

Modern Security Operations Centres (SOCs) are responsible for monitoring increasingly complex network infrastructures spanning LAN, WAN, wireless, cloud, and hybrid environments. The task is usually performed with a combination of firewalls, signature-based IDS, and log aggregation platforms. These tools require continuous manual tuning: new attack signatures must be written, false positives must be triaged, and firewall rules must be maintained. The workload is heavy, the response time is slow, and analysts frequently suffer from alert fatigue. The consequence is that early-stage reconnaissance activity---in particular port scanning---is often noticed only after it has escalated into a more serious intrusion.

\subsection{Problematic}

The core problem is therefore twofold: port scanning and network reconnaissance are the precursors of most targeted attacks, yet current defences are either too slow to react, too rigid to detect novel techniques, or unable to respond autonomously. The situation can be summarised by the following operational question:

\begin{quote}
\textit{Comment peut-on d\'{e}tecter et classifier les attaques de type PortScan \`{a} partir de donn\'{e}es de trafic r\'{e}seau (pcap/csv) afin de prot\'{e}ger l'infrastructure et alerter les \'{e}quipes SOC en temps r\'{e}el?}
\end{quote}

In English: \textit{How can we detect and classify PortScan attacks from network traffic data (pcap/csv) in order to protect the infrastructure and alert SOC teams in real time?} Existing solutions either detect without responding, leaving analysts to act manually, or respond with fixed rules that cannot adapt to stealthy or low-rate scans.

\subsection{Solution Approach}

AEGIS Entropy proposes an AI-driven response that combines two orientations:

\begin{enumerate}
    \item \textbf{Decision-support orientation:} A real-time SOC dashboard that displays classification confidence scores, active scanners, honeypot contacts, and blocked IPs, enabling analysts to understand and validate system decisions.
    \item \textbf{Automation orientation:} An active defence engine that automatically redirects attackers to decoy ports, mutates the exposed attack surface, and blacklists confirmed malicious source addresses without requiring human intervention.
\end{enumerate}

Both orientations are built on the same machine learning core: a set of classifiers trained on network flow features to distinguish benign traffic from port scan activity.

\section{Project Objectives}

The project objectives are defined according to the SMART criteria (Specific, Measurable, Achievable, Relevant, Time-bound). They are summarised in Table~\ref{tab:objectives}.

\begin{table}[H]
\centering
\caption{SMART project objectives}
\label{tab:objectives}
\begin{tabularx}{\textwidth}{lX}
\toprule
\textbf{Objective} & \textbf{Description} \\
\midrule
Detection & Train and deploy an ML model capable of classifying port scan traffic with F1-Score $\geq$ 88\%, Accuracy $\geq$ 90\%, and FPR $<$ 6\% on the CICIDS2017 test subset. \\
Response & Integrate an active defence layer (MTD port mutation, honeypot redirection, automatic blacklisting) triggered by ML detections. \\
Validation & Demonstrate end-to-end operation in a controlled VirtualBox lab with Kali Linux attack scenarios and measure detection/response latency. \\
Usability & Provide a Flask-based SOC dashboard that updates within 5 seconds and presents alerts with confidence scores and severity levels. \\
\bottomrule
\end{tabularx}
\end{table}

\section{Target Audience and User Interviews}\label{sec:interviews}

\subsection{Interview Methodology}

Following the professor's ``Golden Rule'' that a problem must be validated by real users, four semi-structured interviews were conducted with students and future network/security practitioners. The interview guide combined qualitative questions (understanding current difficulties and expectations) and quantitative questions (acceptability thresholds, desired features, fears). The complete transcripts are archived in \texttt{docs/reports/01\_conception/interviews/Guide\_et\_Entretiens\_Complets.md}.

\subsection{Interviewees}

\begin{itemize}
    \item \textbf{Ouzizi Amine} --- Second-year engineering student, Networks and Telecommunications, cybersecurity option.
    \item \textbf{Fahd} --- First-year engineering student, Computer Science.
    \item \textbf{Cybers\'{e}curit\'{e} student} --- Networks and Telecommunications student specialised in cybersecurity.
    \item \textbf{Informatique student} --- Second-year Computer Science student.
\end{itemize}

\subsection{Synthesis of Findings}

Despite different technical backgrounds, the interviewees converged on several key needs and concerns:

\begin{itemize}
    \item \textbf{Manual monitoring is unsustainable.} Interviewees described relying on firewall logs, manual packet captures, or simple ping tests to diagnose suspicious behaviour. They cannot continuously watch traffic while doing their primary work.
    \item \textbf{False positives are a major pain point.} Fahd in particular emphasised that overly strict rules block legitimate academic traffic. Any AI-based system must keep false positives low and provide a clear way to contest automated blocks.
    \item \textbf{Stealthy scans are hard to catch.} The cybersecurity student noted that slow SYN scans and distributed scans easily pass under fixed thresholds. Behavioural analysis over time windows is preferred over instantaneous packet rules.
    \item \textbf{Dynamic defence is welcomed.} Ouzizi and Fahd both viewed dynamic port changes and automated blocking positively, provided the system is transparent and lightweight.
    \item \textbf{Resource consumption and transparency matter.} Interviewees asked for a lightweight web dashboard, clear notifications explaining why a flow was blocked, and guarantees that the AI does not consume excessive CPU resources or take uncontrollable actions.
\end{itemize}

These findings directly shaped the AEGIS Entropy requirements: low FPR, real-time dashboard, explainable decisions, dual detection modes (supervised + unsupervised), and graduated automated response.

\section{Data Collection Strategy}

\subsection{Training Data}

The primary training dataset is the CICIDS2017 PortScan subset~\cite{cicids2017}. It was selected because it is a modern, publicly available benchmark that contains realistic benign and attack traffic at the flow level. The specific file used is \texttt{Friday-WorkingHours-Afternoon-PortScan.pcap\_ISCX.csv}, which contains 286,467 flow records and 79 raw features.

\subsection{Live Capture Data}

In addition to the static training dataset, live traffic is generated in a controlled VirtualBox laboratory. The attacker VM (Kali Linux) runs Nmap scans against the target VM (Ubuntu Server), and the resulting packets and scan logs are used to validate the integration pipeline and dashboard behaviour.

\section{Measurement Instruments and Calibration}\label{sec:instruments}

\subsection{Instruments}

Because the project mixes offline benchmark data with self-collected live traffic, the measurement instruments must be explicitly identified:

\begin{itemize}
    \item \textbf{Nmap:} Used to generate controlled port scan anomalies (SYN, TCP Connect, UDP, XMAS, ACK, service detection, OS detection).
    \item \textbf{Scapy and Wireshark:} Used to capture packets and inspect flow-level behaviour during live demonstrations.
    \item \textbf{CICIDS2017 CSV:} The benchmark instrument providing labelled flow records for training and evaluation.
\end{itemize}

\subsection{Calibration}

Before official collection, the instruments were calibrated and tested in the controlled lab environment. Nmap scan parameters were documented, Scapy capture filters were verified against known scan signatures, and the CICIDS2017 file was checked for completeness, column consistency, and class balance. This calibration step ensures that the data used for training and the data collected during demonstrations are both accurate and reproducible.

\section{Data Dictionary}\label{sec:data_dictionary}

The machine learning model operates on an 11-feature input vector. Nine features are extracted directly from the CICIDS2017 CSV; two are placeholders reserved for real-time integration with the deception subsystem. Table~\ref{tab:data_dict} presents the complete dictionary.

\begin{table}[H]
\centering
\caption{Data dictionary --- input features and target variable}
\label{tab:data_dict}
\begin{tabularx}{\textwidth}{lclX}
\toprule
\textbf{Feature} & \textbf{Type} & \textbf{Source} & \textbf{Description} \\
\midrule
Destination Port & Quantitative & CSV & Target port number; proxy for distinct destination ports \\
Flow Duration & Quantitative & CSV & Duration of the flow in microseconds \\
Total Fwd Packets & Quantitative & CSV & Number of packets sent in the forward direction \\
SYN Flag Count & Quantitative & CSV & Number of SYN-flagged packets \\
RST Flag Count & Quantitative & CSV & Number of RST-flagged packets \\
ACK Flag Count & Quantitative & CSV & Number of ACK-flagged packets \\
Flow IAT Mean & Quantitative & CSV & Mean inter-arrival time between packets \\
Bwd Packet Length Mean & Quantitative & CSV & Mean size of packets in the backward direction \\
Init\_Win\_bytes\_forward & Quantitative & CSV & Initial TCP window size in forward direction \\
MTD Port Delta & Quantitative & Placeholder & Difference between probed port and active MTD port \\
Shadow Node Interaction & Binary & Placeholder & 1 if source contacted a honeypot/decoy, 0 otherwise \\
\midrule
\textbf{Target: Label} & Qualitative & CSV & \textsc{Benign} or \textsc{PortScan} \\
\bottomrule
\end{tabularx}
\end{table}

In offline training, the two placeholder features are initialised to zero. In live deployment they are populated by the deception subsystem described in Chapter~\ref{ch:defense}.

\section{Project Characterization}

\subsection{Type of Artificial Intelligence}

AEGIS Entropy is a \textbf{symbolic-numeric hybrid AI system}. The numeric core consists of machine learning classifiers trained on quantitative flow features; the symbolic layer consists of rules that translate classifier outputs into concrete defence actions (blacklist, redirect, mutate ports).

\subsection{Type of Learning}

The system uses both:

\begin{itemize}
    \item \textbf{Supervised learning} (Random Forest and XGBoost) for binary classification of known scan patterns.
    \item \textbf{Unsupervised learning} (Isolation Forest) for anomaly detection of slow or unknown scan behaviour.
\end{itemize}

\subsection{Type of Model}

The primary models are ensemble tree-based classifiers:

\begin{itemize}
    \item \textbf{Random Forest}~\cite{breiman2001}: bagged decision trees used as the primary detection model.
    \item \textbf{XGBoost}~\cite{chen2016}: gradient-boosted trees used as a comparison benchmark.
    \item \textbf{Isolation Forest}~\cite{liu2008}: unsupervised anomaly detector used as a complementary layer.
\end{itemize}

\subsection{Variables Explicatives and Variable Cible}

The explanatory variables ($X$) are the 11 features listed in Table~\ref{tab:data_dict}. All are numeric; the target variable ($Y$) is the qualitative binary label \textsc{Benign} / \textsc{PortScan}. Class imbalance is handled with SMOTE~\cite{chawla2002} on the training set only, as detailed in Chapter~\ref{ch:data_prep}.

\section{Solution Architecture Overview}

The AEGIS Entropy architecture is organised into three layers:

\begin{itemize}
    \item \textbf{Capture layer:} Network packets are captured with Scapy and aggregated into flow-level feature vectors. Offline mode uses the CICIDS2017 CSV directly; live mode uses the VirtualBox lab.
    \item \textbf{Detection layer:} Features are scaled and passed to the trained Random Forest, XGBoost, and Isolation Forest models. The output is a classification plus a confidence score.
    \item \textbf{Response layer:} Confirmed detections trigger active defence actions: honeypot redirection, MTD port mutation, and source IP blacklisting. Events are streamed to the Flask dashboard via Socket.IO.
\end{itemize}

A detailed description of the capture environment is given in Chapter~\ref{ch:capture}, the defence subsystem in Chapter~\ref{ch:defense}, and the integration layer in Chapter~\ref{ch:integration}.

\section{Conclusion}

This chapter established the problem foundation using the CPS methodology, validated it through four user interviews, and defined the data dictionary that drives the machine learning pipeline. The next chapter describes how the CICIDS2017 dataset is cleaned, explored, and transformed into the feature set used to train the models.
